\section{Idea 7: repeated copies remove the bounded error}\label{sec:copies}
\begin{definition}[Ordered concatenation]\label{def:concat}
Take disjoint label blocks. Put every label of an earlier block above
every label of a later one, and choose an order from the given family
inside each block independently. For integers \(t\geq1\), write \(K^{[t]}\) for \(t\) copies
of one family \(K\).
\end{definition}

\begin{lemma}[Exact product counts]\label{lem:copies}
If \(n,t\geq1\), \(K\in\F_n\), \(M=|K|>0\), then \(K^{[t]}\in\F_{tn}\), and
\[
 |K^{[t]}|=M^t,\quad
 A(K^{[t]})=tA(K)M^{t-1},\quad
 \Sig(K^{[t]})=t\Sig(K)M^{t-1}.
\]
\end{lemma}
\begin{proof}
Retain all internal defining constraints and fix every comparison
between blocks. A triple meeting more than one block has a label
fixed first or last, so a compatible extreme law can be chosen.
The constraints describe exactly the concatenated orders, proving
membership in \(\F\). Independent choices give the first formula.
An enabled internal swap can occur in any of the \(t\) blocks while
the other \(t-1\) blocks have \(M^{t-1}\) choices. A swap across a
join violates a fixed cross-block comparison and is never enabled.
This proves the second formula. Deleting label \(x\) in a block
similarly gives \(M_xM^{t-1}\) outputs. Their blocks are recoverable
from their disjoint labels and forced ordering, so there is no extra
identification or correction at a join. Sum over all blocks and labels.
\end{proof}
\lean{\path{CDCert.K6.concat}, \path{isDom_concat},
\path{card_concat}, \path{adj_concat} formalise the construction and
first two counts at list level. \path{dbl} and \path{dbl_props}
iterate doubling, using \(t=2^k\). The displayed deletion product
is also a written counting proof; the formal amplification uses the
adjacency product instead, then the fibre identity. We do not label
the deletion-product formula as a separately exported Lean theorem.}

For instance \(K=\{12,21\}\) gives
\(K^{[2]}=\{1234,1243,2134,2143\}\). Each deletion leaves two
orders, so \(\Sig=8=2\cdot2\cdot2\); the swap at the join is absent.

\subsection{A slightly larger construction that can be inspected}
Use the four-order domain \(D=\{123,213,231,321\}\) from
Section~\ref{sec:deletion}, and a relabelled copy
\(D'=\{456,546,564,654\}\). Joining them means choosing a prefix
from \(D\) and a suffix from \(D'\). There are 16 orders, including
\(123456,123546,213564,321654\). Every deletion has eight distinct
outputs: two choices inside its affected block, and four in the other.
Thus \(\Sigma=6\cdot8=48\). The two internal path graphs contribute
24 undirected edges, hence \(A=48\). There is no swap across the join.
All fibres are intervals, so \(G=0\) and \(48+0+48=6\cdot16\).

The same recipe gives the following genuine domains. Write
\(n=3t+r\), \(r\in\{0,1,2\}\); join \(t\) copies of \(D\), then
append one fixed label if \(r=1\), or a freely ordered pair if \(r=2\).
Each row counts this constructed family, without claiming it is largest.
\begin{center}\small
\renewcommand{\arraystretch}{1.15}
\begin{tabular}{@{}rrrrrr@{}}\toprule
\(n\)&Three-label blocks&\(M\)&\(\Sigma\)&\(G\)&\(A\)\\\midrule
5&1&8&20&0&20\\
6&2&16&48&0&48\\
7&2&16&64&0&48\\
8&2&32&128&0&128\\
9&3&64&288&0&288\\
10&3&64&352&0&288\\
\bottomrule\end{tabular}\end{center}

This table makes the absence of a new join correction visible in
actual counts. The construction and its counts are checked by the
separate example script; the general product and interval identities
used above are the formalised ones already identified.

\subsection{Remove the outside error}

\begin{theorem}[Exact extreme-law density]\label{thm:density}
For \(n\geq1\) and \(K\in\F_n\),
\[
 \boxed{8A(K)\leq5n|K|,\qquad3n|K|\leq8\Sig(K).}
\]
\end{theorem}
\begin{proof}
Apply Proposition~\ref{prop:affine} once to the whole joined family,
for \(t\) large enough that \(tn\geq7\):
\[
 tA(K)M^{t-1}\leq\frac58tnM^t+BM^t.
\]
Divide by \(tM^{t-1}>0\), obtaining
\(A(K)\leq(5n/8+B/t)M\). Taking arbitrarily large \(t\), or just
powers of two, gives the first inequality. Interval fibres and
\(\Sig+A=nM\) give the second. The original \(n\) can be small,
since only the repeated family's length must reach seven.
\end{proof}
\lean{\path{CDCert.K6.adj_le} proves the first inequality at list
level. Lean avoids a real limit: if the integer bound failed by at
least one, choose \(2^k\) larger than the fixed error can permit,
use \path{dbl_props}, and obtain a contradiction. After transferring
lists back to actual orders and using gap zero,
\path{CDCert.Bridge.InF.density} states the second inequality on
\(\mathrm{Fin}(k+1)\). Both are among the 20 audited results.}
